\documentclass[11pt]{article}
\usepackage[margin=1in]{geometry}
\usepackage{amsmath,amssymb,amsthm,mathtools}
\usepackage{hyperref}

\theoremstyle{plain}
\newtheorem{theorem}{Theorem}
\newtheorem{lemma}[theorem]{Lemma}
\newtheorem{proposition}[theorem]{Proposition}
\newtheorem{corollary}[theorem]{Corollary}
\theoremstyle{definition}
\newtheorem{definition}[theorem]{Definition}
\newtheorem{remark}[theorem]{Remark}
\newtheorem{conjecture}[theorem]{Conjecture}

\DeclareMathOperator{\End}{End}
\newcommand{\F}{\mathcal{F}}
\newcommand{\Om}{\Omega}
\newcommand{\Ombar}{\Om_{\mathrm{bar}}}
\newcommand{\qq}{(q-q^{-1})}
\newcommand{\Bm}{B_{-1}^{[e]}}
\newcommand{\Ce}{C_e^{(1)}}
\newcommand{\N}[1]{N_e^{(#1)}}
\newcommand{\wt}{\operatorname{wt}}

\title{Transposition is the twist: the correction term $\Ce$ as a
height-reversal defect on level-one $q$-Fock space}
\author{Clio}
\date{30 August 2026}

\begin{document}
\maketitle

\begin{abstract}
Let $\F$ be the level-one $q$-deformed Fock space of
$U_q(\widehat{\mathfrak{sl}}_e)$, let $P_e$ be the ribbon-sign operator
$P_e\lvert\lambda\rangle=\sum(-q)^{h}\lvert\mu\rangle$ (sum over
$e$-ribbon additions of height $h$), and let $\Bm$ be the Heisenberg mode
$\sum(-q^{-1})^h$, so that $P_e=\Bm+\qq\Ce$ (result C4).  We prove that
transposition $\Om\lvert\lambda\rangle=\lvert\lambda'\rangle$ \emph{reverses
the ribbon height grading}, $\Om\N{h}\Om=\N{e-1-h}$, and deduce
$\Om P_e\Om=(-q)^{e-1}\Bm$.  Since $\Bm$ is the coefficientwise bar-conjugate
of $P_e$, this identifies the abstract bar involution on this operator with a
combinatorial involution, and answers Q55: $\qq\Ce$ is exactly the deviation
of $P_e$ from its scalar-corrected transposed conjugate.  \textbf{This last
statement is not new}: it is the one-ribbon case of Leclerc--Thibon
Prop.~7.10, which we locate and reconcile scalar-for-scalar in \S\ref{sec:LT};
what is new here is the graded form, and what follows.  We then use the
same circle of ideas to give a \emph{structural} proof --- independent of C4
--- that $[e_i,P_e]$ is divisible by $\qq$, namely $[\bar e_i,P_e]=0$; and we
compute $\Om e_i\Om=q\,K_{-i}^{-1}\,\bar e_{-i}$ (again Leclerc--Thibon
Prop.~7.10, rederived), which shows why $\Om$ is not a crystal morphism.  Sharpening this, we reduce the observed \emph{unit
monomial} form of the quotient, $[e_i,P_e]_{\nu\lambda}=\pm q^m\qq$, to a
purely combinatorial statement about $[e_i,\N{h}]$ --- all the $q$-arithmetic
being forced by the single vanishing $[e_i,\Bm]=0$.  Finally we determine the $i$-signature word of the
horizontal inflation $e\lambda$ to be $A^{\delta}(RA)^m$, which simultaneously
(i) reproves $\varepsilon_i(e\lambda)\le1$ (C5 at level one), (ii) proves
$\varepsilon_i=0$ for the vertical inflation, thereby explaining the
vertical/horizontal dichotomy as prefix-maximum versus suffix-maximum of one
and the same $\pm1$ walk, and (iii) closes the gap in step (3b) of the C5
write-up identified by Lyra.  \textbf{Everything here is level one.}
\end{abstract}

\section{Setting and conventions}

Fix $e\ge2$.  Let $\F=\bigoplus_{\lambda}\mathbb{Q}(q)\lvert\lambda\rangle$,
the sum over all partitions.  We use English notation; the node in row $r$ and
column $c$ (both $1$-indexed) has \emph{residue} $c-r \bmod e$.  For a
partition $\lambda$ we set $\lambda_r=0$ for $r>\ell(\lambda)$ and
$\lambda_0=+\infty$.

A \emph{ribbon} (rim hook) is a finite skew shape $\mu/\lambda$ that is
edge-connected and contains no $2\times2$ square.  For an $e$-ribbon
$\mu/\lambda$ we write $r(\mu/\lambda)$ for the number of rows it meets,
$c(\mu/\lambda)$ for the number of columns, and
\[
  h(\mu/\lambda) := r(\mu/\lambda)-1
\]
for its \emph{height} (equivalently spin/leg length; this convention is fixed
throughout the programme and is the one implemented in
\texttt{probes/2026-08-12-q34-jacon-lacabanne/qfock.py}).

\begin{definition}\label{def:ops}
For $0\le h\le e-1$ let $\N{h}\in\End(\F)$ be the $\mathbb{Q}(q)$-linear
\emph{height-$h$ ribbon counting operator},
\[
  \N{h}\lvert\lambda\rangle
   \;=\; \sum_{\substack{\mu\supset\lambda \\ \mu/\lambda \text{ an $e$-ribbon},\ h(\mu/\lambda)=h}}
          \lvert\mu\rangle .
\]
For an indeterminate $t$ put $R_e(t):=\sum_{h=0}^{e-1}t^h\,\N{h}$.  Then
\[
  P_e = R_e(-q),\qquad \Bm = R_e(-q^{-1}),\qquad
  \Ce=\sum_{h=0}^{e-1}(-1)^h[h]_q\,\N{h},
\]
where $[h]_q=(q^h-q^{-h})/(q-q^{-1})$.  Result~C4 is the identity
$P_e=\Bm+\qq\Ce$, i.e.\ $(-q)^h-(-q^{-1})^h=(-1)^h\qq[h]_q$.
\end{definition}

Let $\mathrm{b}\colon\End(\F)\to\End(\F)$ be the coefficientwise bar
involution: apply $q\mapsto q^{-1}$ to every matrix entry in the basis
$\{\lvert\lambda\rangle\}$.  It is a ring automorphism over $\mathbb{Z}$
(semilinear over $\mathbb{Z}[q^{\pm1}]$).  We write $\bar X=\mathrm{b}(X)$.
Since the $\N{h}$ have integer entries,
\begin{equation}\label{eq:barP}
  \overline{P_e}=\Bm,\qquad \overline{\Ce}=\Ce
\end{equation}
(the second because $[h]_q$ is bar-invariant).

Let $\Om\in\End(\F)$ be transposition, $\Om\lvert\lambda\rangle=
\lvert\lambda'\rangle$, extended $\mathbb{Q}(q)$-linearly; $\Om^2=\mathrm{id}$.
Let $\Ombar$ be the $q$-\emph{antilinear} transposition,
$\Ombar(\sum a_\lambda\lvert\lambda\rangle)=\sum\overline{a_\lambda}
\lvert\lambda'\rangle$; also an involution.  Write
$\mathrm{t}(X):=\Om X\Om$; then $\Ombar X\,\Ombar=\mathrm{t}(\bar X)$.

\section{Transposition reverses the height grading}

\begin{lemma}[Ribbon rectangle]\label{lem:rc}
If $\mu/\lambda$ is an $e$-ribbon then $r(\mu/\lambda)+c(\mu/\lambda)=e+1$.
Equivalently $c=e-h$ where $h=r-1$.
\end{lemma}

\begin{proof}
Let $S=\mu/\lambda$ and let $r$ be the number of rows it meets.  In each such
row $S$ occupies a nonempty contiguous interval of columns; relabel the rows
of $S$ as $1,\dots,r$ from top to bottom and let row $i$ occupy $[a_i,b_i]$.
Because $\lambda$ and $\mu$ are partitions, $a_i\ge a_{i+1}$ and
$b_i\ge b_{i+1}$.  Edge-connectivity of $S$ forces consecutive rows to share
at least one column, i.e.\ $a_i\le b_{i+1}$; absence of a $2\times2$ square
forces them to share at most one, and their overlap is exactly
$[a_i,b_{i+1}]$.  Hence
\[
  b_{i+1}=a_i \qquad (1\le i\le r-1).
\]
In particular the column intervals chain together, so the set of columns met
by $S$ is exactly $[a_r,b_1]$ and $c=b_1-a_r+1$.  Now
\[
  e=\sum_{i=1}^{r}(b_i-a_i+1)
   =\sum_{i=1}^{r}b_i-\sum_{i=1}^{r}a_i+r
   =\sum_{i=1}^{r}b_i-\Bigl(\sum_{i=2}^{r}b_i+a_r\Bigr)+r
   =b_1-a_r+r=(c-1)+r . \qedhere
\]
\end{proof}

\begin{lemma}[Transposition of ribbon additions]\label{lem:transpose}
For partitions $\lambda\subset\mu$, the shape $\mu/\lambda$ is an $e$-ribbon
if and only if $\mu'/\lambda'$ is, and in that case
$r(\mu'/\lambda')=c(\mu/\lambda)$ and $c(\mu'/\lambda')=r(\mu/\lambda)$.
\end{lemma}

\begin{proof}
The map $(r,c)\mapsto(c,r)$ is a bijection of $\mathbb{Z}_{>0}^2$ carrying the
Young diagram of $\nu$ to that of $\nu'$ for every partition $\nu$; it
preserves inclusion, cardinality and edge-adjacency (hence connectivity),
carries $2\times2$ squares to $2\times2$ squares, and interchanges rows with
columns.
\end{proof}

\begin{theorem}[Height reversal]\label{thm:main}
For all $e\ge2$ and $0\le h\le e-1$,
\[
  \boxed{\;\Om\,\N{h}\,\Om \;=\; \N{e-1-h}\;}
\]
\end{theorem}

\begin{proof}
Both sides are $\mathbb{Q}(q)$-linear, so it suffices to evaluate on
$\lvert\lambda\rangle$.  Since $\Om^2=\mathrm{id}$,
\[
  \Om\N{h}\Om\lvert\lambda\rangle
  =\Om\N{h}\lvert\lambda'\rangle
  =\Om\!\!\sum_{\substack{\nu/\lambda' \text{ $e$-ribbon}\\ h(\nu/\lambda')=h}}\!\!\lvert\nu\rangle
  =\!\!\sum_{\substack{\nu/\lambda' \text{ $e$-ribbon}\\ r(\nu/\lambda')=h+1}}\!\!\lvert\nu'\rangle .
\]
By Lemma~\ref{lem:transpose} the assignment $\nu\mapsto\nu'$ is a bijection
from $\{\nu:\nu/\lambda' \text{ an $e$-ribbon with } r=h+1\}$ onto
$\{\mu:\mu/\lambda \text{ an $e$-ribbon with } c=h+1\}$, and by
Lemma~\ref{lem:rc} the latter condition reads $r(\mu/\lambda)=e+1-(h+1)=e-h$,
i.e.\ $h(\mu/\lambda)=e-1-h$.  The right-hand sum is therefore exactly
$\N{e-1-h}\lvert\lambda\rangle$.
\end{proof}

\begin{corollary}[Functional equation]\label{cor:funeq}
$\Om\,R_e(t)\,\Om=t^{\,e-1}R_e(t^{-1})$ as an identity of operators over
$\mathbb{Z}[t^{\pm1}]$.
\end{corollary}
\begin{proof}
$\Om R_e(t)\Om=\sum_h t^h\N{e-1-h}=\sum_k t^{e-1-k}\N{k}=t^{e-1}R_e(t^{-1})$.
\end{proof}

\begin{corollary}[The boxed identity of the brief]\label{cor:boxed}
\[
  \Om\,P_e\,\Om=(-q)^{e-1}\,\Bm,
  \qquad
  \Om\,\Bm\,\Om=(-q^{-1})^{e-1}\,P_e .
\]
\end{corollary}
\begin{proof}
Put $t=-q$ resp.\ $t=-q^{-1}$ in Corollary~\ref{cor:funeq} and use
$(-q)^{-1}=-q^{-1}$.
\end{proof}

\begin{corollary}[Transposition \emph{is} bar, on this operator]\label{cor:tb}
$\mathrm{t}(P_e)=(-q)^{e-1}\,\mathrm{b}(P_e)$ and
$\mathrm{t}(\Bm)=(-q^{-1})^{e-1}\,\mathrm{b}(\Bm)$.  Equivalently, with
$\Ombar$ the antilinear transposition,
\[
  \Ombar\,P_e\,\Ombar=(-q^{-1})^{e-1}P_e ,
\]
so $P_e$ is an eigenvector of the involution $X\mapsto\Ombar X\Ombar$ with
eigenvalue $(-q^{-1})^{e-1}$.
\end{corollary}
\begin{proof}
Combine Corollary~\ref{cor:boxed} with \eqref{eq:barP}.
\end{proof}

\begin{corollary}[The spin normalisation trivialises the scalar]\label{cor:spin}
Work over $\mathcal{A}=\mathbb{Z}[t^{\pm1/2}]$ with $t=-q$ and the involution
$t\mapsto t^{-1}$ (which restricts to $q\mapsto q^{-1}$ on
$\mathbb{Z}[q^{\pm1}]$, since $t^{-1}=-q^{-1}$).  Put
$s=t^{(e-1)/2}$ and $\widehat P_e:=s^{-1}P_e$.  Then $\bar s=s^{-1}$ and
\[
  \Ombar\,\widehat P_e\,\Ombar=\widehat P_e .
\]
\end{corollary}
\begin{proof}
$\Ombar(cX)\Ombar=\bar c\,\Ombar X\Ombar$ for a scalar $c$, so by
Corollary~\ref{cor:tb}, $\Ombar\widehat P_e\Ombar=
\bar s^{-1}\cdot t^{-(e-1)}P_e=s\,s^{-2}P_e=s^{-1}P_e=\widehat P_e$,
using $\bar s=t^{-(e-1)/2}=s^{-1}$ and $t^{-(e-1)}=(-q^{-1})^{e-1}$.
\end{proof}

\begin{remark}
Corollary~\ref{cor:spin} is a Hilbert-90-shaped statement: the scalar
$c=(-q^{-1})^{e-1}$ satisfies $c\,\bar c=1$, and the ``spin shift'' by
$(e-1)/2$ --- the standard centring of the LLT spin statistic --- is
precisely a solution of $c=\bar s/s$.  The half-integer power is unavoidable
for $e$ even, exactly as in LLT theory.
\end{remark}

\section{Relation to Leclerc--Thibon Proposition 7.10}
\label{sec:LT}

\textbf{Corollaries~\ref{cor:boxed}--\ref{cor:spin} are not new.}  They are
the one-ribbon case of a published result, which I located on disk only after
proving them, and I record the dictionary here rather than leave the
attribution to a reader.

Lam (\texttt{math/0310250}, \verb|lam.tex:1026-1041|, read directly)
introduces the \emph{semi-linear} involution $v\mapsto v'$ on $\F$ defined by
$q'=q^{-1}$ and $\lvert\lambda\rangle\mapsto\lvert\lambda'\rangle$ --- which
is exactly our $\Ombar$ --- and quotes:

\begin{proposition}[Leclerc--Thibon, ASPM \textbf{28} (2000), Prop.~7.10;
quoted at \texttt{lam.tex:1030}]\label{prop:LT710}
For all $u\in\F$ and compositions $\beta$ with $|\beta|=k$,
\[
  (e_iu)'=q^{\,h_{-i}-1}e_{-i}u',\qquad
  (f_iu)'=q^{-h_{-i}-1}f_{-i}u',
\]
\[
  (\mathcal{V}_\beta u)'=(-q)^{(n-1)k}\widetilde{\mathcal{V}}_\beta u',
  \qquad
  (\mathcal{U}_\beta u)'=(-q)^{(n-1)k}\widetilde{\mathcal{U}}_\beta u',
\]
where $\mathcal{V},\mathcal{U}$ are the horizontal and
$\widetilde{\mathcal{V}},\widetilde{\mathcal{U}}$ the vertical
$n$-ribbon-strip operators.
\end{proposition}

The dictionary, checked rather than assumed:
\begin{itemize}
\item $n=e$.
\item Lam's spin of a single ribbon is its height, $\text{rows}-1$
  (\verb|lam.tex:429-436|: ``the height of the ribbon $r$ will then be called
  its spin $s(r)$''), i.e.\ our $h$.
\item $\mathcal{V}_k\lvert\lambda\rangle=\sum_\mu(-q)^{-s(\mu/\lambda)}
  \lvert\mu\rangle$ (\verb|lam.tex:848|).  At $k=1$ a horizontal ribbon strip
  of size $1$ is a single $e$-ribbon and $(-q)^{-h}=(-q^{-1})^{h}$, so
  \[
     \mathcal{V}_1=\Bm .
  \]
  A single ribbon is also a vertical strip of size $1$, so
  $\widetilde{\mathcal{V}}_1=\mathcal{V}_1=\Bm$ as well.
\item Hence Prop.~\ref{prop:LT710} at $k=1$ reads
  $\Ombar\,\Bm\,\Ombar=(-q)^{e-1}\Bm$.  Our Corollary~\ref{cor:boxed} gives
  $\Ombar\Bm\Ombar=\Om\overline{\Bm}\Om=\Om P_e\Om=(-q)^{e-1}\Bm$.
  \textbf{The two agree, including the scalar.}
\item Our Proposition~\ref{prop:omega-chev} is the first two relations of
  Prop.~\ref{prop:LT710}: barring $\Om e_i\Om=q K_{-i}^{-1}\bar e_{-i}$ gives
  $\Ombar e_i\Ombar=q^{-1}K_{-i}e_{-i}$, which is $(e_iu)'=q^{h_{-i}-1}e_{-i}u'$.
\end{itemize}

Two open questions from the 30 August reading log are therefore settled.
\emph{Q56 --- is the linear $\Om$ the same map as LT's semilinear involution?}
No: they are different maps, related by $\Ombar=\mathrm{b}\circ\Om$, and both
statements are true in their own form (Corollary~\ref{cor:boxed} for $\Om$,
Corollary~\ref{cor:tb} for $\Ombar$); the passage between them is exactly
$\Bm=\overline{P_e}$.  \emph{Q57 --- does spin/cospin exchange account for the
exponent?}  Yes: Lemma~\ref{lem:rc} ($r+c=e+1$) \emph{is} the spin/cospin
exchange for a single ribbon, and at $k=1$ the exponent $(n-1)k=e-1$ matches
on the nose.

What is not subsumed:
\begin{itemize}
\item Theorem~\ref{thm:main} is a \emph{graded} statement, $q$-free, at the
  level of the individual $\N{h}$; Prop.~\ref{prop:LT710} is its
  $q$-specialised shadow at $k=1$.  (For general $k$ the graded analogue is
  the statement that conjugation exchanges spin and cospin on ribbon tableaux,
  $s(T')=\mathrm{mspin}(sh(T))-s(T)$, which Lam also discusses.)
\item Conversely Prop.~\ref{prop:LT710} at $k\ge2$ is \emph{stronger} than
  anything proved here: it treats strips of several ribbons, where horizontal
  and vertical strips genuinely differ.  We prove only $k=1$.
\item Theorem~\ref{thm:t2}, Proposition~\ref{prop:reduction},
  Conjecture~\ref{conj:sharp}, Theorem~\ref{thm:word},
  Corollary~\ref{cor:dichotomy} and the closing of the C5 step-(3b) gap are
  not consequences of Prop.~\ref{prop:LT710} and are, as far as I know, new.
\end{itemize}

\section{What transposition does to $\Ce$}

\begin{proposition}\label{prop:Cim}
\[
  \Om\,\Ce\,\Om \;=\; (-1)^{e-1}\sum_{k=0}^{e-1}(-1)^{k}\,[\,e-1-k\,]_q\;\N{k}.
\]
In particular, for $e\ge2$, $\Om\Ce\Om$ is neither a scalar multiple of $\Ce$
nor equal to $\overline{\Ce}$.
\end{proposition}
\begin{proof}
Apply Theorem~\ref{thm:main} termwise to $\Ce=\sum_h(-1)^h[h]_q\N{h}$ and
reindex $k=e-1-h$, using $(-1)^{h}=(-1)^{e-1-k}$.
For the last claim take $e=2$: $\Ce=-\N{1}$ while $\Om\Ce\Om=-\N{0}$, and
$\N{0},\N{1}$ are linearly independent (already on $\lvert\varnothing\rangle$,
where $\N{0}\lvert\varnothing\rangle=\lvert(2)\rangle$ and
$\N{1}\lvert\varnothing\rangle=\lvert(1,1)\rangle$).  Also
$\overline{\Ce}=\Ce$ by \eqref{eq:barP}.
\end{proof}

So the correction term is \emph{not} transposition-(anti)symmetric.  Under
$\Om$ the quantum integer $[h]_q$ attached to height $h$ is replaced by the
\emph{complementary} quantum integer $[e-1-h]_q$.  This was one of the three
candidates listed in the brief; the other two are refuted by
Proposition~\ref{prop:Cim} (and, computationally, by $381$ explicit failures
each --- see \S\ref{sec:tally}).

\section{The answer to Q55}

\begin{theorem}\label{thm:q55}
Define the $\mathbb{Z}[q^{\pm1}]$-linear map
$\delta_\Om\colon\End(\F)\to\End(\F)$,
$\delta_\Om(X):=X-(-q)^{-(e-1)}\,\Om X\Om$.  Then
\[
  \delta_\Om(P_e)\;=\;P_e-\Bm\;=\;\qq\,\Ce .
\]
\end{theorem}
\begin{proof}
Immediate from Corollary~\ref{cor:boxed} and C4.
\end{proof}

Thus $\qq\Ce$ \emph{is} a coboundary, for the following action and for no
other reason: the order-two twisted conjugation $X\mapsto(-q)^{-(e-1)}\Om X\Om$
on $\End(\F)$, which is an honest involution after the spin renormalisation of
Corollary~\ref{cor:spin} (before it, its square is multiplication by
$(-q)^{-2(e-1)}$).  We state it that way deliberately: this is not a statement
of group cohomology.

Two further remarks, both of which limit the headline.

\begin{corollary}\label{cor:notartifact}
$\Ce=0$ if and only if $e=1$.  In particular $\Ce$ is \emph{not} an artifact
of a normalisation: no rescaling $P_e\mapsto uP_e$ makes it vanish.
\end{corollary}
\begin{proof}
$\Ce=\sum_h(-1)^h[h]_q\N{h}$; the $\N{h}$ for $0\le h\le e-1$ are nonzero and
linearly independent (evaluate on $\lvert\varnothing\rangle$, where the
$e$-ribbon additions are exactly the hooks, one of each height), and
$[h]_q\ne0$ for $h\ge1$.
\end{proof}

\begin{remark}
What \emph{is} an artifact of a convention is the \emph{scalar} $(-q)^{e-1}$
in Corollary~\ref{cor:boxed}: it disappears in the spin-centred normalisation.
The correction term itself measures a genuine reversal asymmetry, namely that
$h\mapsto[h]_q$ is not symmetric under $h\mapsto e-1-h$.  The brief's
optimistic reading (``$\Ce$ demoted from mysterious correction to artifact of
an inflation convention'') is therefore \emph{too strong} and is not claimed
here.  What is achieved is the identification of the abstract bar involution
$\mathrm{b}$, on this operator, with the concrete combinatorial involution
$\Om$.
\end{remark}

\section{A structural proof of the $\qq$-divisibility of $[e_i,P_e]$}
\label{sec:t2}

We now fix the Leclerc--Thibon action of $U_q(\widehat{\mathfrak{sl}}_e)$ on
$\F$.  For $\gamma$ an addable or removable $i$-node of $\lambda$ set
\[
  N^{>}_i(\lambda,\gamma)=\#\{\text{addable $i$-nodes of $\lambda$ strictly
  above }\gamma\}-\#\{\text{removable $i$-nodes strictly above}\},
\]
and $N^{<}_i$ likewise with ``below''.  Then
\[
  f_i\lvert\lambda\rangle=\sum_{\gamma \text{ addable }i}
     q^{\,N^{>}_i(\lambda,\gamma)}\lvert\lambda+\gamma\rangle,
  \qquad
  e_i\lvert\lambda\rangle=\sum_{\gamma \text{ removable }i}
     q^{-N^{<}_i(\lambda\!-\!\gamma,\gamma)}\lvert\lambda-\gamma\rangle,
\]
and $K_i\lvert\lambda\rangle=q^{\langle h_i,\wt\lambda\rangle}
\lvert\lambda\rangle$ with $\langle h_i,\wt\lambda\rangle=
A_i(\lambda)-R_i(\lambda)$ (numbers of addable and removable $i$-nodes).  We
verified $[e_i,f_j]=\delta_{ij}(K_i-K_i^{-1})/\qq$ on $665$ instances as an
independent witness that this is the genuine module (\S\ref{sec:tally}).

We take as given (Uglov; used already in C4) the level-one commutation
\begin{equation}\label{eq:uglov}
  [e_i,\Bm]=[f_i,\Bm]=0 \qquad (i\in\mathbb{Z}/e).
\end{equation}

\begin{theorem}\label{thm:t2}
Let $\bar e_i=\mathrm{b}(e_i)$, $\bar f_i=\mathrm{b}(f_i)$ and put
\[
  E_i:=\frac{e_i-\bar e_i}{\qq},\qquad F_i:=\frac{f_i-\bar f_i}{\qq}.
\]
Then $E_i,F_i$ have entries in $\mathbb{Z}[q^{\pm1}]$ --- explicitly, the
entry of $E_i$ at $(\lambda-\gamma,\lambda)$ is $-[\,N^{<}_i\,]_q$ and that of
$F_i$ at $(\lambda+\gamma,\lambda)$ is $[\,N^{>}_i\,]_q$ --- and
\[
  [\bar e_i,P_e]=[\bar f_i,P_e]=0,
  \qquad
  [e_i,P_e]=\qq\,[E_i,P_e],
  \qquad
  [f_i,P_e]=\qq\,[F_i,P_e].
\]
In particular every entry of $[e_i,P_e]$ and of $[f_i,P_e]$ is divisible by
$\qq$ in $\mathbb{Z}[q^{\pm1}]$.  (This is Conjecture~4 of
\texttt{2026-08-12-commutation-defect.tex} and the statement the registry
calls ``T2''; see Remark~\ref{rem:labels} for the label collision.)
\end{theorem}

\begin{proof}
$\mathrm{b}$ is a ring automorphism of $\End(\F)$ over $\mathbb{Z}$, so
applying it to \eqref{eq:uglov} and using $\mathrm{b}(\Bm)=P_e$ from
\eqref{eq:barP} gives $[\bar e_i,P_e]=[\bar f_i,P_e]=0$.  Subtracting this
from $[e_i,P_e]$ gives $[e_i,P_e]=[e_i-\bar e_i,P_e]$.  Finally $e_i-\bar e_i$
is entrywise $q^{-N}-q^{N}=-\qq[N]_q$, hence divisible by $\qq$ with quotient
in $\mathbb{Z}[q^{\pm1}]$; likewise $f_i-\bar f_i$ is entrywise
$q^{N}-q^{-N}=\qq[N]_q$.  (Here $[-N]_q=-[N]_q$, so negative $N$ is covered.)
\end{proof}

This is the structural reason the brief asked for, and it is
\emph{independent of C4}: it does not use the decomposition
$P_e=\Bm+\qq\Ce$, only \eqref{eq:uglov} and the fact that $P_e=\overline{\Bm}$.
The slogan is: \emph{$\Bm$ commutes with the Chevalley operators, so $P_e$
commutes with their bar-conjugates, and the two differ by $\qq$ times an
integral operator.}  Comparing with the C4 route gives a nontrivial identity.

\begin{corollary}\label{cor:tworoutes}
$[e_i,\Ce]=[E_i,P_e]$ and $[f_i,\Ce]=[F_i,P_e]$.
\end{corollary}
\begin{proof}
By C4 and \eqref{eq:uglov}, $[e_i,P_e]=\qq[e_i,\Ce]$; compare with
Theorem~\ref{thm:t2} and cancel $\qq$ (which is a nonzerodivisor).
\end{proof}

\begin{remark}[A label collision, resolved]\label{rem:labels}
The names ``T2/T3'' are used differently in two places in this programme, and
the difference matters.  In \texttt{2026-08-12-commutation-defect.tex}, T2 is
the $q=1$ statement $[e_i,P_e]|_{q=1}=0$ and T3 is the $(q-1)$-divisibility of
$e_i\cdot v_{k',e}$; the $\qq$-divisibility is \emph{Conjecture~4} there,
explicitly flagged as empirical.  In the registry node
\texttt{T2-qdiff-divisibility} and in the 30~August brief, ``T2'' names the
$\qq$-divisibility of $[e_i,P_e]$ and ``T3'' the claim that the quotient is
linear in $q$.  Theorem~\ref{thm:t2} above proves the \emph{mathematical}
statement --- $\qq$-divisibility of every entry of $[e_i,P_e]$ and $[f_i,P_e]$
--- which was the 2026-08-12 Conjecture~4 and is the registry's ``T2''.  It
does not touch the 2026-08-12 T2/T3, which are proved there already.  No
result is affected; the labels are.
\end{remark}

\subsection{A sharper divisibility, and a clean reduction}
\label{sec:sharp}

The registry's ``T3'' --- that the quotient is linear in $q$ --- is not what
the data say.  The data say something much stronger and much simpler.

\begin{proposition}\label{prop:reduction}
Fix $e\ge2$, $i\in\mathbb{Z}/e$, $x\in\{e_i,f_i\}$ and partitions
$\lambda,\nu$.  Put $c_h:=\bigl([x,\N{h}]\bigr)_{\nu\lambda}\in
\mathbb{Z}[q^{\pm1}]$ and $H:=\{h: c_h\ne0\}$.  Then:
\begin{enumerate}
\item $|H|\ne1$.
\item Suppose $|H|=2$, say $H=\{b,\,b+d\}$ with $d\ge1$, and suppose
  $c_b=\varepsilon q^{k}$ with $\varepsilon\in\{\pm1\}$.  Then necessarily
  $c_{b+d}=-\varepsilon(-1)^{d}q^{k+d}$, and
  \[
    \bigl([x,P_e]\bigr)_{\nu\lambda}
      \;=\; -\varepsilon(-1)^{b}\,q^{\,k+b+d}\;\qq\;[d]_q .
  \]
  In particular the quotient by $\qq$ is a unit monomial $\pm q^{m}$ if and
  only if $d=1$.
\end{enumerate}
\end{proposition}

\begin{proof}
Let $\Phi(t):=\sum_{h=0}^{e-1}c_h t^h\in\mathbb{Z}[q^{\pm1}][t]$.  Since
$R_e(t)$ is $\mathbb{Z}[t]$-linear in the $\N{h}$, we have
$\Phi(t)=([x,R_e(t)])_{\nu\lambda}$; hence
\[
  \Phi(-q^{-1})=([x,\Bm])_{\nu\lambda}=0
  \qquad\text{and}\qquad
  \Phi(-q)=([x,P_e])_{\nu\lambda},
\]
the first by \eqref{eq:uglov}.

(1) If $H=\{b\}$ then $0=\Phi(-q^{-1})=c_b(-q^{-1})^{b}$, so $c_b=0$, a
contradiction.

(2) From $c_{b+d}(-q^{-1})^{b+d}+c_b(-q^{-1})^{b}=0$ we get
$c_{b+d}=-c_b(-1)^{d}q^{d}=-\varepsilon(-1)^{d}q^{k+d}$.  Then
\[
  \Phi(-q)=c_{b+d}(-q)^{b+d}+c_b(-q)^{b}
   =\varepsilon(-1)^{b}\Bigl[-q^{k+d+b+d}+q^{k+b}\Bigr]
   =-\varepsilon(-1)^{b}q^{k+b}\bigl(q^{2d}-1\bigr),
\]
using $(-1)^{d}(-1)^{b+d}=(-1)^{b}$.  Finally
$q^{2d}-1=q^{d}\bigl(q^{d}-q^{-d}\bigr)=q^{d}\qq[d]_q$.
\end{proof}

So \emph{all} the $q$-arithmetic is forced by the single vanishing
$\Phi(-q^{-1})=0$.  What remains is purely combinatorial.

\begin{conjecture}\label{conj:sharp}
For all $e\ge2$, $i\in\mathbb{Z}/e$, $x\in\{e_i,f_i\}$ and all $\lambda,\nu$:
\begin{enumerate}
\item[\textup{(P2)}] every entry $\bigl([x,\N{h}]\bigr)_{\nu\lambda}$ is
  either $0$ or a unit monomial $\pm q^{k}$; and
\item[\textup{(P1)}] the support $H$ is either empty or a pair of
  \emph{consecutive} integers.
\end{enumerate}
By Proposition~\ref{prop:reduction} these imply that every entry of
$[e_i,P_e]$ and of $[f_i,P_e]$ is exactly $\pm q^{m}\qq$ for some
$m\in\mathbb{Z}$ --- a unit monomial quotient, far stronger than linearity in
$q$.
\end{conjecture}

Evidence, over $e\le6$ and $|\lambda|\le10$: (P2) holds for all $5366$
nonzero entries of the operators $[x,\N{h}]$ tested; (P1) holds for all
$2683$ nonempty supports, every one of which has $|H|=2$ with the two indices
consecutive; and, computed directly and independently of the reduction, all
$2683$ nonzero entries of $[e_i,P_e]$ and $[f_i,P_e]$ are of the form
$\pm q^{m}\qq$.  Zero exceptions in each case.

Note that Proposition~\ref{prop:reduction}(1) already rules out $|H|=1$
unconditionally, so the content of (P1) is the bound $|H|\le2$ together with
consecutiveness; and $|H|=0$ is exactly the case of a vanishing commutator
entry.  Theorem~\ref{thm:t2} remains the only \emph{proved} divisibility
statement here; Conjecture~\ref{conj:sharp} is the sharp form, and the
combinatorial statement it reduces to --- an explicit commutation rule
between a Chevalley operator and ribbon addition at fixed height --- looks to
me like the right next target, replacing the ill-posed ``T3 linearity''.

\section{Transposition on the Chevalley generators}

\begin{proposition}\label{prop:omega-chev}
For all $i\in\mathbb{Z}/e$,
\[
  \Om\,e_i\,\Om \;=\; q\,K_{-i}^{-1}\,\bar e_{-i},
  \qquad
  \Om\,f_i\,\Om \;=\; \bar f_{-i}\,\cdot\,q^{-1}K_{-i},
\]
where in the first identity $K_{-i}^{-1}$ acts on the target and in the second
$K_{-i}$ acts on the source.
\end{proposition}

\begin{proof}
Transposition sends the node $(r,c)$ to $(c,r)$, hence sends residue $i$ to
$-i$, addable nodes to addable nodes, removable to removable, and the relation
``strictly above'' to ``strictly to the left''.  By Lemma~\ref{lem:order}
below, within the set of addable and removable $i$-nodes of a partition,
``strictly to the left'' coincides with ``strictly below''.  Hence
$N^{<}_i(\lambda,\gamma)=N^{>}_{-i}(\lambda',\gamma')$ and
$N^{>}_i(\lambda,\gamma)=N^{<}_{-i}(\lambda',\gamma')$.  Therefore
$\Om e_i\Om$ removes $(-i)$-nodes with coefficient $q^{-N^{>}_{-i}}$ and
$\Om f_i\Om$ adds $(-i)$-nodes with coefficient $q^{N^{<}_{-i}}$.  Now for a
node $\gamma$ addable to $\nu$,
\[
  N^{>}_j(\nu,\gamma)+N^{<}_j(\nu,\gamma)
  = \bigl(A_j(\nu)-1\bigr)-R_j(\nu)
  = \langle h_j,\wt\nu\rangle-1 ,
\]
since $\gamma$ itself is one of the $A_j(\nu)$ addable $j$-nodes.  Applying
this with $\nu=\lambda-\gamma$ (for $e_i$) and $\nu=\lambda$ (for $f_i$)
converts $q^{-N^{>}}$ into $q^{\,1-\langle h\rangle}q^{N^{<}}$ and $q^{N^{<}}$
into $q^{\langle h\rangle-1}q^{-N^{>}}$, which are the two stated formulas.
\end{proof}

\begin{corollary}\label{cor:t2-transposed}
$[\,\Om e_i\Om,\;P_e\,]=[\,\Om f_i\Om,\;P_e\,]=0$.
\end{corollary}
\begin{proof}
Conjugate \eqref{eq:uglov} by $\Om$ and use Corollary~\ref{cor:boxed}:
$0=\Om[e_i,\Bm]\Om=[\Om e_i\Om,(-q^{-1})^{e-1}P_e]$.
\end{proof}

So $P_e$ commutes with two visibly different deformations of the Chevalley
action --- the bar-conjugate one and the transposed one --- and
Proposition~\ref{prop:omega-chev} says they differ only by the diagonal
operator $q^{\pm1}K_{\mp i}^{\mp1}$, with which $P_e$ commutes because an
$e$-ribbon contains exactly one node of each residue and so preserves
$\langle h_j,\wt\cdot\rangle$.  The two statements are one statement.

\section{The crystal side: why $\Om$ is not a crystal morphism}
\label{sec:crystal}

\begin{lemma}[Order reversal]\label{lem:order}
Let $e\ge2$, $\lambda$ a partition, $i\in\mathbb{Z}/e$, and let
$\mathcal{N}_i(\lambda)$ be the set of addable and removable $i$-nodes of
$\lambda$.  If $(r_1,c_1),(r_2,c_2)\in\mathcal{N}_i(\lambda)$ with $r_1<r_2$
then $c_1>c_2$.
\end{lemma}

\begin{proof}
First note that every $(r,c)\in\mathcal{N}_i(\lambda)$ satisfies
\begin{equation}\label{eq:node}
  \lambda_r\le c\le\lambda_r+1,\qquad \lambda_{r-1}\ge c,\qquad
  \lambda_{r+1}\le c-1 .
\end{equation}
Indeed a removable node is $(r,\lambda_r)$ with $\lambda_{r+1}<\lambda_r$, and
an addable node is $(r,\lambda_r+1)$ with $\lambda_{r-1}>\lambda_r$, i.e.\
$\lambda_{r-1}\ge\lambda_r+1$.

Now take $r_1<r_2$.  If $r_2\ge r_1+2$ then
$\lambda_{r_2}\le\lambda_{r_1+1}\le c_1-1$ by \eqref{eq:node}, so
$c_2\le\lambda_{r_2}+1\le c_1$.  If equality held we would have
$c_2=c_1$ and $\lambda_{r_2}=c_2-1$, so $(r_2,c_2)$ is addable and
$\lambda_{r_2-1}\ge c_2=c_1$; but $r_2-1\ge r_1+1$ gives
$\lambda_{r_2-1}\le\lambda_{r_1+1}\le c_1-1$, a contradiction.  Hence
$c_1>c_2$.

If $r_2=r_1+1$ then again $\lambda_{r_2}\le c_1-1$ and $c_2\le c_1$, and
equality forces $\lambda_{r_2}=c_1-1$, $(r_2,c_2)$ addable, and
$\lambda_{r_1}=\lambda_{r_2-1}\ge c_1$; combined with $\lambda_{r_1}\le c_1$
this gives $\lambda_{r_1}=c_1$, so $(r_1,c_1)$ is removable.  The two residues
are then $c_1-r_1$ and $c_1-r_1-1$, which are distinct modulo $e$ because
$e\ge2$ --- contradicting that both nodes have residue $i$.  Hence $c_1>c_2$.
\end{proof}

\begin{definition}\label{def:word}
The \emph{$i$-signature word} $w_i(\lambda)\in\{A,R\}^*$ is obtained by
listing the elements of $\mathcal{N}_i(\lambda)$ in order of increasing row
and recording $A$ for addable, $R$ for removable.  For $w\in\{A,R\}^*$ let
$\mathrm{exc}(u)=\#R(u)-\#A(u)$ and set
\[
  \varepsilon(w)=\max_{u\ \text{prefix of}\ w}\mathrm{exc}(u),
  \qquad
  \varepsilon^{\mathrm{op}}(w)=\max_{u\ \text{suffix of}\ w}\mathrm{exc}(u)
\]
(both maxima include the empty word, hence are $\ge0$).  Then
$\varepsilon_i(\lambda)=\varepsilon(w_i(\lambda))$ is the Kashiwara
$\varepsilon$ of the level-one Fock crystal: it is the number of $R$'s
surviving repeated cancellation of factors $AR$.  Dually
$\varepsilon^{\mathrm{op}}$ counts the $R$'s surviving cancellation of factors
$RA$.
\end{definition}

\begin{theorem}\label{thm:eps-transpose}
For every partition $\lambda$ and every $i$,
\[
  w_{-i}(\lambda')=\overleftarrow{w_i(\lambda)}
  \qquad\text{and hence}\qquad
  \varepsilon_{-i}(\lambda')=\varepsilon^{\mathrm{op}}\bigl(w_i(\lambda)\bigr).
\]
\end{theorem}
\begin{proof}
Transposition is a bijection $\mathcal{N}_i(\lambda)\to
\mathcal{N}_{-i}(\lambda')$ preserving the type ($A$ or $R$) and exchanging
rows with columns.  By Lemma~\ref{lem:order} the row order on
$\mathcal{N}_i(\lambda)$ is the reverse of the column order, so reading
$\mathcal{N}_{-i}(\lambda')$ by increasing row means reading
$\mathcal{N}_i(\lambda)$ by increasing column, i.e.\ by decreasing row.  For
the second statement, reversing a word exchanges prefixes with suffixes.
\end{proof}

Since $\varepsilon\ne\varepsilon^{\mathrm{op}}$ in general (already on
$w=RA$, where $\varepsilon=1$ and $\varepsilon^{\mathrm{op}}=0$), $\Om$ is not
a morphism of $\widehat{\mathfrak{sl}}_e$-crystals; we recorded $490$ explicit
instances of $\varepsilon_{-i}(\lambda')\ne\varepsilon_i(\lambda)$ in the
range tested.  This is consistent with, and explained by,
Proposition~\ref{prop:omega-chev}: $\Om e_i\Om$ is not $e_{-i}$ but its
bar-conjugate twisted by a diagonal, and bar exchanges the lower crystal
lattice with the upper one.

\subsection{The signature word of a horizontal inflation}

Write $e\lambda:=(e\lambda_1,e\lambda_2,\dots)$ for the \emph{horizontal}
inflation, and $V(\lambda)$ for the \emph{vertical} inflation, in which every
part of $\lambda$ is repeated $e$ times.  These are exchanged by
transposition:
\begin{equation}\label{eq:HV}
  V(\lambda)=\bigl(e\lambda'\bigr)' .
\end{equation}
(The transpose of $e\mu$ has each part of $\mu'$ repeated $e$ times.)

\begin{theorem}\label{thm:word}
Let $\lambda$ be a partition, $\ell=\ell(\lambda)$, $\mu=e\lambda$, and
$i\in\mathbb{Z}/e$.  Then
\[
  w_i(\mu)=A^{\delta}\,(RA)^{m},
  \qquad
  \delta=[\,i\equiv0\,],\quad
  m=\#\{\,1\le r\le\ell \;:\; r\equiv-i \bmod e,\ \lambda_r>\lambda_{r+1}\,\}.
\]
\end{theorem}

\begin{proof}
Since $\mu_r=e\lambda_r\equiv0\pmod e$:
\begin{itemize}
\item the removable nodes of $\mu$ are the $(r,e\lambda_r)$ with
  $1\le r\le\ell$ and $\lambda_r>\lambda_{r+1}$; such a node has residue
  $e\lambda_r-r\equiv-r$;
\item the addable nodes of $\mu$ are the $(s,e\lambda_s+1)$ with
  $1\le s\le\ell+1$ and $\lambda_{s-1}>\lambda_s$ (recall
  $\lambda_0=+\infty$, $\lambda_{\ell+1}=0$); such a node has residue
  $e\lambda_s+1-s\equiv1-s$.
\end{itemize}
Substituting $s=r+1$ turns the addable condition into
$\lambda_{r}>\lambda_{r+1}$ and the addable residue into $-r$.  Hence:
\begin{equation}\label{eq:Sr}
  \boxed{\begin{aligned}
  &\text{$\mu$ has a removable node in row $r$}
     &&\iff&& \lambda_r>\lambda_{r+1},\\
  &\text{$\mu$ has an addable node in row $r+1$}
     &&\iff&& \lambda_r>\lambda_{r+1},
  \end{aligned}}
\end{equation}
and when they exist both have residue $-r \bmod e$.  Consequently
$\mathcal{N}_i(\mu)$ is the disjoint union of
\begin{itemize}
\item the pairs $\{(r,e\lambda_r),\,(r+1,e\lambda_{r+1}+1)\}$ over
  $1\le r\le\ell$ with $r\equiv-i$ and $\lambda_r>\lambda_{r+1}$; and
\item the single addable node $(1,1)$, present iff $1\equiv1-i$, i.e.\ iff
  $i\equiv0$ (this is the only addable node whose partner index $r=0$ lies
  outside $[1,\ell]$; note $\lambda_0>\lambda_1$ holds by convention).
\end{itemize}
Each pair occupies the consecutive rows $r,r+1$ and contributes $RA$ in
reading order.  Distinct pairs have starting rows differing by a multiple of
$e\ge2$, so they occupy disjoint and non-interleaving row blocks
$\{r,r+1\}$, $\{r+e,r+e+1\}$, \dots  The lone addable node, when present, is
in row $1$, while pairs then start at rows $\equiv0\bmod e$, i.e.\ at
$r\ge e\ge2$; so it precedes them all.  This is exactly $A^\delta(RA)^m$.
\end{proof}

\begin{corollary}\label{cor:dichotomy}
For all $\lambda$ and all $i$:
\begin{enumerate}
\item $\varepsilon_i(e\lambda)\le1$, with equality iff $i\not\equiv0$ and
  $m\ge1$; in particular $\varepsilon_0(e\lambda)=0$ always.  \emph{(This is
  C5 at level one, reproved.)}
\item $\varepsilon^{\mathrm{op}}\bigl(w_i(e\lambda)\bigr)=0$.
\item $\varepsilon_i\bigl(V(\lambda)\bigr)=0$.
\end{enumerate}
\end{corollary}
\begin{proof}
By Theorem~\ref{thm:word} the word is $A^\delta(RA)^m$.  Its prefixes have
excess $\le1$, with $1$ attained iff the word begins with $R$, i.e.\ $\delta=0$
and $m\ge1$; this gives (1).  Its suffixes all end in $A$ (if nonempty) and
have excess $\le0$ --- the suffix $(RA)^k$ has excess $0$ and
$A(RA)^k$ has excess $-1$ --- which gives (2).  For (3), combine
\eqref{eq:HV} with Theorem~\ref{thm:eps-transpose}:
$\varepsilon_i(V(\lambda))=\varepsilon_i\bigl((e\lambda')'\bigr)
=\varepsilon^{\mathrm{op}}\bigl(w_{-i}(e\lambda')\bigr)=0$.
\end{proof}

This is the promised reconciliation.  \emph{The vertical/horizontal dichotomy
is prefix-maximum versus suffix-maximum of one and the same $\pm1$ walk.}  The
word $RARA\cdots RA$ has a lone uncancelled $R$ at its front and none at its
back; transposition reverses it; the operator statement
(Corollary~\ref{cor:boxed}) and the crystal statement
(Corollary~\ref{cor:dichotomy}) are two readings of the same involution $\Om$.

\subsection{A gap in the C5 write-up, closed}

Lyra's review of \texttt{2026-08-11-C5-gerber-bicrystal.tex} (30 August 2026)
observed that step (3b) asserts, without proof, that the addable node of row
$r$ and the removable node of row $r-1$ both exist precisely when
$\lambda_{r-1}>\lambda_{r}$, and identified this coincidence of booleans as
the load-bearing fact.  Display \eqref{eq:Sr} in the proof of
Theorem~\ref{thm:word} is exactly that statement, and its proof is the single
observation that every row of $e\lambda$ has length divisible by $e$: the
corner of row $r$ then sits in a column $\equiv0$, giving residue $-r$, while
the addable node of row $r+1$ sits in a column $\equiv1$, giving residue
$-r$ as well; and both validity conditions unwind to $\lambda_r>\lambda_{r+1}$
because inflation by $e$ is strictly increasing on parts.  The gap is closed.

\section{Computational verification}
\label{sec:tally}

All probes are in \texttt{probes/2026-08-30-Q55-transpose/}.  Ribbon additions
were computed by two independent routes --- an abacus/$\beta$-number route
(the implementation used throughout this programme) and a brute-force
Young-diagram route that enumerates all $\mu\supset\lambda$ with $\mu/\lambda$
connected, of size $e$, and free of $2\times2$ squares.  The two agree on all
$239$ pairs $(e,\lambda)$ tested, which pins the height convention
$h=\text{rows}-1$ in both.

\begin{center}
\begin{tabular}{llrr}
\hline
Statement & Range & Checks & Failures\\
\hline
route A $=$ route B (ribbon additions) & $e\le6$, $|\lambda|\le8$ & $239$ & $0$\\
Lemma~\ref{lem:rc} \ $r+c=e+1$ & $e\le6$, $|\lambda|\le8$ & $1052$ & $0$\\
Thm~\ref{thm:main} (partition level) & $e\le8$, $|\lambda|\le14$ & $3556$ ($21910$ additions) & $0$\\
Thm~\ref{thm:main} (operator level, per $h$) & $e\le7$, $|\lambda|\le9$ & $1458$ & $0$\\
Cor.~\ref{cor:boxed} \ $\Om P_e\Om=(-q)^{e-1}\Bm$ & $e\le7$, $|\lambda|\le9$ & $381$ & $0$\\
Cor.~\ref{cor:boxed} \ $\Om \Bm\Om=(-q^{-1})^{e-1}P_e$ & $e\le7$, $|\lambda|\le9$ & $381$ & $0$\\
Cor.~\ref{cor:tb} \ $\Ombar P_e\Ombar=(-q^{-1})^{e-1}P_e$ & $e\le7$, $|\lambda|\le9$ & $381$ & $0$\\
Prop.~\ref{prop:Cim} (complementary form) & $e\le7$, $|\lambda|\le9$ & $381$ & $0$\\
\ \ \emph{refuted}: $\Om\Ce\Om=-(-q)^{e-1}\Ce$ & $e\le7$, $|\lambda|\le9$ & $381$ & $381$\\
\ \ \emph{refuted}: $\Om\Ce\Om=\overline{\Ce}$ & $e\le7$, $|\lambda|\le9$ & $381$ & $381$\\
\hline
$[e_i,f_i]=(K_i-K_i^{-1})/\qq$ & $e\le5$, $|\lambda|\le8$ & $665$ & $0$\\
\eqref{eq:uglov} \ $[e_i,\Bm]=[f_i,\Bm]=0$ & $e\le5$, $|\lambda|\le8$ & $2\times665$ & $0$\\
Thm~\ref{thm:t2} \ $[\bar e_i,P_e]=[\bar f_i,P_e]=0$ & $e\le5$, $|\lambda|\le8$ & $2\times665$ & $0$\\
Thm~\ref{thm:t2} \ $[e_i,P_e]=\qq[E_i,P_e]$ & $e\le5$, $|\lambda|\le8$ & $665$ & $0$\\
Thm~\ref{thm:t2} \ $[f_i,P_e]=\qq[F_i,P_e]$ & $e\le5$, $|\lambda|\le8$ & $665$ & $0$\\
Cor.~\ref{cor:tworoutes} \ $[e_i,\Ce]=[E_i,P_e]$ & $e\le5$, $|\lambda|\le8$ & $665$ & $0$\\
Cor.~\ref{cor:tworoutes} \ $[f_i,\Ce]=[F_i,P_e]$ & $e\le5$, $|\lambda|\le8$ & $665$ & $0$\\
Prop.~\ref{prop:omega-chev} (both forms, $e$ and $f$) & $e\le5$, $|\lambda|\le8$ & $4\times665$ & $0$\\
\hline
Lemma~\ref{lem:order} (columns decrease) & $e\le6$, $|\lambda|\le10$ & $2780$ & $0$\\
Thm~\ref{thm:eps-transpose} (word reversal) & $e\le6$, $|\lambda|\le10$ & $2780$ & $0$\\
Thm~\ref{thm:eps-transpose} ($\varepsilon$ form) & $e\le6$, $|\lambda|\le10$ & $2780$ & $0$\\
\eqref{eq:HV} \ $V(\lambda)=(e\lambda')'$ & $e\le6$, $|\lambda|\le8$ & $335$ & $0$\\
Cor.~\ref{cor:dichotomy}(1) $\varepsilon_i(e\lambda)\le1$ & $e\le6$, $|\lambda|\le8$ & $1340$ & $0$\\
Cor.~\ref{cor:dichotomy}(2) $\varepsilon^{\mathrm{op}}=0$ & $e\le6$, $|\lambda|\le8$ & $1340$ & $0$\\
Cor.~\ref{cor:dichotomy}(3) $\varepsilon_i(V(\lambda))=0$ & $e\le6$, $|\lambda|\le8$ & $1340$ & $0$\\
\eqref{eq:Sr} \ the $S_r$ coupling (both halves) & $e\le7$, $|\lambda|\le10$ & $2\times4890$ & $0$\\
\eqref{eq:Sr} \ residues both $\equiv-r$ & $e\le7$, $|\lambda|\le10$ & $2\times1704$ & $0$\\
Thm~\ref{thm:word} \ $w_i(e\lambda)=A^\delta(RA)^m$ & $e\le7$, $|\lambda|\le10$ & $3753$ & $0$\\
\hline
Conj.~\ref{conj:sharp}(P2) unit monomial entries & $e\le6$, $|\lambda|\le10$ & $5366$ & $0$\\
Conj.~\ref{conj:sharp}(P1) $|H|=2$, consecutive & $e\le6$, $|\lambda|\le10$ & $2683$ & $0$\\
$[x_i,P_e]_{\nu\lambda}=\pm q^m\qq$ (direct) & $e\le6$, $|\lambda|\le10$ & $2683$ & $0$\\
\hline
\end{tabular}
\end{center}

Non-vacuity: in the range $e\le6$, $|\lambda|\le10$ there are $490$ instances
with $\varepsilon_{-i}(\lambda')\ne\varepsilon_i(\lambda)$, so
Theorem~\ref{thm:eps-transpose} is not asserting an equality that holds for
trivial reasons.  The two \emph{refuted} rows above are the other two
candidates the brief asked to be tested; they fail on every instance where the
commutator is nonzero, the smallest being $e=2$, $\lambda=\varnothing$.

\section{What is not closed}
\label{sec:open}

\begin{enumerate}
\item \textbf{The sharp divisibility is not proved.}
Theorem~\ref{thm:t2} gives $\qq$-divisibility structurally, but not the
observed \emph{unit monomial} quotient.  Proposition~\ref{prop:reduction}
reduces that to the purely combinatorial Conjecture~\ref{conj:sharp}
((P1)+(P2)), which is the natural next target.  Note that the route through
$E_i$ in Theorem~\ref{thm:t2} cannot see it: the entries of $E_i$ are quantum
integers $[N^{<}_i]_q$ of unbounded degree, so the monomial shape must arise
from cancellation inside $[E_i,P_e]$, which that argument does not control.
The registry's phrasing ``the quotient is linear in $q$'' should be retired in
favour of the monomial statement (see Remark~\ref{rem:labels}).
\item \textbf{Level $\ge2$.}  Everything above is level one.
Lemma~\ref{lem:rc} and Theorem~\ref{thm:main} are statements about single
ribbons and should lift verbatim to each component of a multipartition; the
identification of $\Om$ with the vertical/horizontal exchange, and
Theorem~\ref{thm:word}, will \emph{not} lift unchanged, because the level-$\ell$
signature is a merge by content of the per-component words and transposition
must also act on the multicharge $s\mapsto-s$ (reversed).  Nothing here
resolves the $\ell\ge2$ obstruction recorded at
\texttt{C5-level-ell-individual-bound}, and nothing here needs Uglov Prop.~3.16.
\item \textbf{$\Ombar$ and the canonical basis --- partly answered, in the
literature.}  $\Ombar$ is exactly the semilinear involution of
Prop.~\ref{prop:LT710} (\S\ref{sec:LT}), and Lam draws from that proposition
the immediate consequence $(G_{n\lambda})'=(-q)^{(n-1)k}G_{n\lambda'}$: so
$\Ombar$ \emph{does} permute the canonical basis elements indexed by
inflations, up to the same scalar.  What is \emph{not} settled is the
behaviour on $G_\mu$ for general $\mu$, and hence whether this is a genuine
transposition symmetry of the LLT polynomials rather than a statement about
the $e$-inflated locus.  That is the question I would take next on this
branch, and it should be attacked by reading LT ASPM~28 \S7 in the original
rather than through Lam's citation.
\item \textbf{The ``artifact'' reading is refuted.}  Corollary~\ref{cor:notartifact}
shows $\Ce\ne0$ for every $e\ge2$ in every normalisation.  The brief's hoped-for
demotion of $\Ce$ to a convention artifact does not survive; what survives is
the identification of the bar involution with transposition on this operator.
\item \textbf{Rick's $\beta'$ programme.}  Nothing here is a claim about
Rick's objects.  The present results are consistent with the suspicion that
his ``vertical versus connected'' obstruction and the inflation-convention
obstruction are the same one --- Corollary~\ref{cor:dichotomy} makes that
precise \emph{at level one, for $\varepsilon_i$}, and no further --- but the
identification of his $\kappa_\mu$ with $P_2$ has not been checked and is not
asserted.
\end{enumerate}

\end{document}
